\scanpage{32}
\begin{center}$-\ 30\ -$\end{center}
\hfill n$^\circ$ 371

\uline{tensoriel de cette famille et on note} $\bigotimes_{i\in I}E_i$ (ou $E_1\otimes_A E_2\otimes\cdots\otimes_A E_n$, \uline{ou} $E_1\otimes E_2\otimes\cdots\otimes E_n$ \uline{lorsque} $I=\left[\!\left[1,n\right]\!\right]$) \uline{l'algèbre (n'ayant pas nécessairement d'élément unité) obtenue en munissant le} $A$-\uline{module} \inserted{produit} \uline{tensoriel des} $A$-\uline{modules} $E_i$ \uline{de la multiplication définie par} (1).

Il est clair que si chacune des $E_i$ admet un élément unité $e_i$, $e=\bigotimes_i e_i$ est élément unité de $E=\bigotimes_{i\in I}E_i$. Si pour chaque $i\in I$, $E_i^\circ$ est l'algèbre \uline{opposée} à $E_i$, le produit tensoriel $\bigotimes_{i\in I}E_i^\circ$ est canoniquement isomorphe à l'algèbre $E^\circ$ opposée à $E$, à laquelle on l'identifie ; en particulier, si les $E_i$ sont \uline{commutatives}, il en est de même de $E$.

Si $(E_i)_{i\in I}$, $(F_i)_{i\in I}$ sont deux familles finies de $A$-algèbres et, pour chaque $i\in I$, $u_i:E_i\to F_i$ un homomorphisme \uline{d'algèbres}, alors
\[
u=\bigotimes_i u_i:\bigotimes_{i\in I}E_i\longrightarrow\bigotimes_{i\in I}F_i
\]
est un homomorphisme d'algèbres, comme il résulte aussitôt de (1).

Pour toute partition $(I_j)_{j\in J}$ de $I$, les isomorphismes d'associativité
\[
\bigotimes_{j\in J}\left(\bigotimes_{i\in I_j}E_i\right)\longrightarrow\bigotimes_{i\in I}E_i
\]
(chap. II, $3^{\mathrm e}$ éd., \S\ 3, n$^\circ$ 9) sont aussi des isomorphismes d'\uline{algèbres}, comme il résulte de leur définition.

Nous nous bornerons dans le reste de ce n$^\circ$ aux propriétés des produits tensoriels de deux algèbres, laissant au lecteur le soin de les étendre aux produits tensoriels de familles finies quelconques.

Soient $E,F$ deux $A$-algèbres ; si $a$ (resp. $b$) est un idéal à gauche de $E$ (resp. $F$), l'image canonique $\operatorname{Im}(a\otimes b)$ de $a\otimes_A b$ dans $E\otimes_A F$ est un idéal à gauche de $E\otimes F$ ; on a des énoncés analogues en remplaçant ``idéal à gauche'' par ``idéal à droite'' ou ``idéal bilatère''. En outre :
\mspage{31}
